% =====================================================================
\section{Vue d'ensemble}

\subsection{La pile technique}

\noindent\begin{tabularx}{\textwidth}{|>{\columncolor{btkpale}\bfseries}L{3.4cm}|X|}
\hline
Couche & Technologie retenue\\
\hline
Présentation & JSP + JSTL, CSS écrit à la main (\texttt{css/style.css}, 794 lignes), graphiques ApexCharts, icônes Lucide\\
\hline
Contrôle & Servlets Jakarta annotées \texttt{@WebServlet}, plus un filtre \texttt{@WebFilter("/*")}\\
\hline
Métier & EJB \texttt{@Stateless} (et un \texttt{@Singleton} pour la sécurité), injectés par \texttt{@EJB}\\
\hline
Persistance & JPA / Hibernate, \texttt{@PersistenceContext}, requêtes JPQL\\
\hline
Transactions & JTA, gérées par le conteneur (CMT) --- aucun \texttt{commit} écrit à la main\\
\hline
Base de données & Oracle, accédée par la datasource \texttt{java:/OracleDS}\\
\hline
Serveur & WildFly, racine de contexte \texttt{/gestion-agences}\\
\hline
Construction & Maven, empaquetage \texttt{war}, JDK 21\\
\hline
Décisionnel & Vues \texttt{V\_BI\_*} en Oracle, chaîne ETL Python, Power BI Desktop\\
\hline
\end{tabularx}

\begin{retenir}
\textbf{En une phrase.} Une application web Jakarta EE 10 en couches, servie
par WildFly, qui lit et écrit dans Oracle par JPA, et dont les mêmes données
alimentent un tableau de bord embarqué et un rapport Power BI par
l'intermédiaire de vues dédiées et d'une chaîne ETL.
\end{retenir}

\subsection{Ce que contient le projet}

Inventaire relevé dans le dépôt :

\noindent\begin{tabularx}{\textwidth}{|>{\columncolor{btkpale}\bfseries}L{3.6cm}|c|X|}
\hline
Élément & Nombre & Emplacement\\
\hline
Servlets & 18 & \texttt{src/main/java/controller/}\\
\hline
Services EJB & 16 & \texttt{src/main/java/service/}\\
\hline
Entités JPA & 16 & \texttt{src/main/java/entity/}\\
\hline
Filtre & 1 & \texttt{src/main/java/filter/NotificationFilter.java}\\
\hline
Pages JSP & 17 & \texttt{src/main/webapp/} et \texttt{WEB-INF/}\\
\hline
Feuille de style & 1 & \texttt{src/main/webapp/css/style.css}\\
\hline
Vues décisionnelles & 4 & \texttt{sql/v\_bi\_*.sql}, \texttt{sql/setup\_pointage.sql}\\
\hline
\end{tabularx}

\subsection{L'architecture en quatre couches}

Chaque requête traverse les mêmes quatre couches, dans le même ordre. C'est le
point que le jury vérifie en premier : \textbf{aucune couche n'en saute une
autre}. Un servlet ne parle jamais à la base ; un service ne fabrique jamais de
HTML.

\begin{center}
\small
\begin{tabular}{@{}c@{}}
\colorbox{btkpale}{\parbox{0.72\textwidth}{\centering\textbf{Navigateur}\\[2pt]\footnotesize requête HTTP}}\\[4pt]
$\downarrow$\\[4pt]
\colorbox{btkpale}{\parbox{0.72\textwidth}{\centering\textbf{1. Présentation} --- \texttt{*.jsp} + \texttt{style.css} + ApexCharts\\[2pt]\footnotesize affiche, ne calcule pas}}\\[4pt]
$\downarrow$\\[4pt]
\colorbox{btkpale}{\parbox{0.72\textwidth}{\centering\textbf{2. Contrôle} --- \texttt{*Servlet.java} (\texttt{@WebServlet})\\[2pt]\footnotesize lit les paramètres, contrôle la session et le rôle, choisit la vue}}\\[4pt]
$\downarrow$\\[4pt]
\colorbox{btkpale}{\parbox{0.72\textwidth}{\centering\textbf{3. Métier} --- \texttt{*Service.java} (\texttt{@Stateless})\\[2pt]\footnotesize règles de gestion, transactions, journalisation, notifications}}\\[4pt]
$\downarrow$\\[4pt]
\colorbox{btkpale}{\parbox{0.72\textwidth}{\centering\textbf{4. Persistance} --- JPA / Hibernate (\texttt{EntityManager})\\[2pt]\footnotesize traduit les objets en SQL Oracle}}\\[4pt]
$\downarrow$\\[4pt]
\colorbox{btkpale}{\parbox{0.72\textwidth}{\centering\textbf{Base Oracle} --- datasource \texttt{java:/OracleDS}}}
\end{tabular}
\end{center}

\subsection{Le cycle de vie d'une requête, sur un exemple}

Prenons l'affichage de la liste des agences, \texttt{/agences}. C'est le
parcours à savoir raconter de mémoire.

\begin{enumerate}[leftmargin=*,itemsep=2pt]
  \item Le navigateur demande \texttt{GET /gestion-agences/agences}.
  \item Le \textbf{filtre} \texttt{NotificationFilter} s'exécute d'abord : il
        lit le rôle en session et pose dans la requête les compteurs de
        demandes en attente, pour que le bandeau les affiche.
  \item WildFly route vers \texttt{AgenceServlet}, associé à \texttt{/agences}
        par l'annotation \texttt{@WebServlet}.
  \item Le servlet vérifie la \textbf{session} et le \textbf{rôle}. Sans
        session valide, redirection vers la page de connexion.
  \item Il appelle \texttt{agenceService.findAll()}. L'EJB ouvre une
        transaction JTA, exécute une requête JPQL, Hibernate la traduit en SQL,
        Oracle répond.
  \item Le servlet range le résultat dans la requête
        (\texttt{setAttribute}) et \textbf{transfère} vers
        \texttt{/WEB-INF/agences.jsp}.
  \item La JSP parcourt la liste avec \texttt{<c:forEach>} et produit le HTML,
        qui part au navigateur.
\end{enumerate}

\begin{retenir}
\textbf{Transfert ou redirection ?} Un \texttt{forward} reste côté serveur :
l'URL ne change pas et les attributs de requête survivent --- c'est ce qu'on
utilise pour \emph{afficher}. Un \texttt{sendRedirect} demande au navigateur de
refaire une requête : l'URL change et les attributs sont perdus --- c'est ce
qu'on utilise \emph{après une écriture}, pour qu'un rafraîchissement ne
rejoue pas l'enregistrement.
\end{retenir}

\subsection{Pourquoi les JSP sont sous \texttt{WEB-INF}}

Toutes les vues sauf deux sont dans \texttt{src/main/webapp/WEB-INF/}. Le
conteneur \textbf{interdit} d'y accéder par une URL directe : on ne peut y
entrer que par un \texttt{forward} depuis un servlet. Conséquence, aucune page
ne peut être ouverte en contournant le contrôle de session et de rôle.

Les deux exceptions sont volontaires : \texttt{index.jsp} (page de connexion)
et \texttt{forgot-password.jsp}, qui doivent être atteignables sans être
authentifié.
